\documentclass[10pt,a4paper]{amsart}
\usepackage[margin=19mm]{geometry}
\usepackage{amsmath,amssymb,mathtools}
\usepackage[scr=rsfs,cal=euler]{mathalfa}
\usepackage{xfrac}
\usepackage[unicode,hidelinks,bookmarks=false]{hyperref}
\newcommand{\ie}{i.e.\ }
\newcommand{\cf}{cf.\ }
\newcommand{\resp}{respectively\ }
\newcommand{\Subsection}{\subsection}
\providecommand{\nobreakdash}{\nobreak\hbox{-}}
\setlength{\emergencystretch}{2em}
\multlinegap0pt
\usepackage{physics}
\theoremstyle{plain}
\newtheorem{conjecture}{Conjecture}
\newtheorem{lemma}{Lemma}
\newtheorem{theorem}{Theorem}
\newtheorem{observation}{Observation}
\newtheorem{definition}{Definition}
\newtheorem{thm}{Theorem}
\newtheorem{lem}{Lemma}
\newtheorem{prop}{Proposition}
\newtheorem{defn}{Definition}
\newtheorem{corollary}{Corollary}
\newtheorem{remark}{Remark}
\title{Stronger Welch Bounds and Optimal Approximate $k$-Designs}
\author{Riccardo Castellano, Dmitry Grinko, Sadra Boreiri, Nicolas Brunner and Jef Pauwels}
\date{Source: arXiv:2602.13099v3 (CC BY 4.0); the authors' own native \LaTeX{} source, set here in the editorial AMS \texttt{amsart} wrapper (the APS \texttt{revtex4-2} class named by the source is not redistributed).}
\begin{document}
\maketitle
\noindent\textit{Editorial scope.} The counted claim of this unit is the article's Proposition 2 (source label \texttt{prop:PT-toric-Haar-spectrum}), the spectrum of the partially transposed Haar moment operator, delivered together with the authors' complete original proof. The statement and the proof are the only retained passages, and every one of their characters is the authors' own; the proof is delivered as the single primary proof body exactly as the authors' proof environment encloses it. Printed numbering: the authors' own declarations of the theorem-like environments are carried unchanged, each environment keeping its own counter as in the authors' preamble, and the proposition counter is advanced so that the delivered PDF prints the counted claim as Proposition~2, exactly as the published article does. Omitted: the rest of the article, that is its title block, abstract, introduction, the Welch-bound and design sections, the main technical results that precede this proposition and the sections that follow it; nothing of the counted statement or of the counted proof is omitted. Class substitution: the source class is the APS \texttt{revtex4-2} class with the authors' own packages, including \texttt{physics} for the Dirac notation used inside the proof; that class is not present in this package and is not redistributed, and the retained author text is set in an editorial AMS \texttt{amsart} wrapper carrying the authors' own macros and environment declarations. Reference handling: no retained passage contains a citation; the article's own TeX contains no reference list and its bibliography exists only in the authors' file \texttt{biblio.bib}, which is kept verbatim in the eprint directory and is not redistributed. No reference entry is transcribed, delivered or invented. No mathematics is written, repaired or re-derived here.

\setcounter{prop}{1}
\begin{prop}[Spectrum of the partially transposed Haar moment]
    \label{prop:PT-toric-Haar-spectrum}
     For $b\in\{0,\dotsc,t\}$, set $a=t-b$ and let $\mathsf{\Gamma}_{b}$ denote the partial transposition over the last $b$ subsystems. The nonzero eigenvalues of $\Omega_t^{\mathsf{\Gamma}_b}$ are
    \begin{equation}
        \lambda_v
        =
        \frac{1}{m^t}
        \sum_{\substack{|\alpha|=a,\ |\beta|=b\\
                        \alpha-\beta=v}}
        \binom{a}{\alpha}
        \binom{b}{\beta},
        \qquad
        v\in\mathcal{V}_{a,b}^{(m)},
        \label{eq:PT-toric-eigenvalues}
    \end{equation}
    Thus,
    \begin{equation}
        \operatorname{rank}\!\left(\Omega_t^{\mathsf{\Gamma}_b}\right)
        =
        \left|\mathcal{V}_{a,b}^{(m)}\right|.
        \label{eq:PT-toric-rank}
    \end{equation}
\end{prop}

\begin{proof}

For every $v\in\mathcal{V}_{a,b}^{(m)}$, define
    \begin{equation}
        \ket{w_v}
        \coloneqq
        \sum_{\substack{|\alpha|=a,\ |\beta|=b\\
                        \alpha-\beta=v}}
        \sqrt{
            \binom{a}{\alpha}
            \binom{b}{\beta}
        }\,
        \ket{\alpha}\otimes\ket{\beta}.
        \label{eq:wv-definition}
    \end{equation}
    We will prove that
    \begin{equation}
        \Omega_t^{\mathsf{\Gamma}_b}
        =
        \frac{1}{m^t}
        \sum_{v\in\mathcal{V}_{a,b}^{(m)}}
        \ketbra{w_v},
        \label{eq:PT-toric-rank-one-decomposition}
    \end{equation}
    and that the vectors $\ket{w_v}$ are mutually orthogonal.

    First, we have
\begin{equation}
    \Omega_{t}^{\mathsf{\Gamma}_{b}}
    =
    \int_{P(\mathbb{T}^{m})}\ketbra{\phi}
    ^{\otimes a}\otimes\ketbra{\overline{\phi}}
    ^{\otimes b}
    \mathrm{d}\mu(\phi).
    \label{eq:toric-Haar-moment2}
\end{equation}


    
    In the normalized occupation-number bases one has
    \begin{align}
        \ket{\phi}^{\otimes a}
        \otimes
        \ket{\overline{\phi}}^{\otimes b}
        &=
        \frac{1}{m^{t/2}}
        \sum_{\substack{|\alpha|=a\\|\beta|=b}}
        \sqrt{
            \binom{a}{\alpha}
            \binom{b}{\beta}
        }\,
        e^{\mathrm{i}(\alpha-\beta)\cdot\phi}
        \ket{\alpha}\otimes\ket{\beta}
        \nonumber\\
        &=
        \frac{1}{m^{t/2}}
        \sum_{v\in\mathcal{V}_{a,b}^{(m)}}
        e^{\mathrm{i}v\cdot\phi}\ket{w_v}.
        \label{eq:occupation-expansion-PT}
    \end{align}
    Since every $v\in\mathcal{V}_{a,b}^{(m)}$ satisfies
    $\sum_jv_j=a-b$, the difference $v-v'$ belongs to the root lattice
    $\mathbb{Z}_{0}^{m}$.
   
    Orthogonality of the characters of $P(\mathbb{T}^{m})$ therefore gives
    \begin{equation}
        \int_{P(\mathbb{T}^{m})}
        e^{\mathrm{i}(v-v')\cdot\phi}
        \,\mathrm{d}\mu([\phi])
        =
        \delta_{v,v'}.
    \end{equation}
    Substituting \eqref{eq:occupation-expansion-PT} into
    \eqref{eq:toric-Haar-moment2} proves
    \eqref{eq:PT-toric-rank-one-decomposition}.

    Finally, if $v\neq v'$, the vectors $\ket{w_v}$ and
    $\ket{w_{v'}}$ are supported on disjoint sets of pairs
    $(\alpha,\beta)$, since a given pair determines a unique difference
    $\alpha-\beta$. They are therefore orthogonal. Every
    $\ket{w_v}$ is non-zero by the definition of
    $\mathcal{V}_{a,b}^{(m)}$, so each summand contributes exactly one
    nonzero eigenvalue, proving the claim.
\end{proof}

\end{document}
